\section{The Northcott theorem}

Throughout this section, $S=(K,(\Omega,\mathcal A,\nu),\phi)$ is a
proper adelic curve and $\overline K$ is an algebraic closure of $K$.
Heights on $\overline K$ are taken with respect to
$S\otimes_K\overline K$ as in Definition~\ref{def:absolute-projective-height}.
We prove the assertions of \cite[Section 3.5]{chen-moriwaki} using the
finite-embedding compatibility of the extension measures. For polynomial
products we retain the archimedean exponent in the estimate; its integral
is finite even when the measure of the archimedean locus is infinite.

\begin{lemma}[Heights under finite field embeddings]
\label{lem:height-finite-embedding}
\uses{def:tuple-height,lem:local-height-integrable,lem:algebraic-finite-embedding-projection}
\lean{ArakelovGeometry.AdelicStructure.localHeight_comp_finiteEmbedding,ArakelovGeometry.AdelicStructure.logHeight_comp_finiteEmbedding,ArakelovGeometry.AdelicStructure.affineHeight_comp_finiteEmbedding}
\leanok
Let $S$ be any adelic structure, let $E/K$ be finite, let $L/K$ be algebraic, and let
$i:E\hookrightarrow L$ be a $K$-embedding. If $I$ is finite and nonempty
and $a=(a_j)_{j\in I}\in E^I$, then
\[
 h_{S\otimes_K L}\bigl((i(a_j))_{j\in I}\bigr)
 =h_{S\otimes_K E}(a).
\]
In particular, the height of an algebraic tuple can be computed over
any finite extension containing its coordinates.
\end{lemma}
\begin{proof}
Let $p_i:\Omega_L\to\Omega_E$ be pullback of absolute values along
$i$. It is measurable, its measure pushforward is $\nu_E$, and
$|i(b)|_x=|b|_{p_i(x)}$. Taking the maximum over the coordinates and
then its logarithm gives
$\ell_{S\otimes_K L}(i(a),x)=\ell_{S\otimes_K E}(a,p_i(x))$.
This also holds for the zero tuple by its stipulated zero local height.
The local height on $\Omega_E$ is measurable and integrable. The
pushforward identity gives equality of the integrals of its nonnegative
positive and negative parts and of its absolute value. Subtracting the
first two equalities proves the displayed equality of heights. Every
finite family of algebraic elements generates a finite extension, and
its inclusion in $\overline K$ gives the last assertion.
\end{proof}

\begin{lemma}[Heights of conjugates]
\label{lem:height-conjugate}
\dcref{Lemma 3.5.4}
\source{chen-moriwaki-adelic:page-0216}
\uses{lem:height-finite-embedding,def:affine-height,def:absolute-projective-height}
\lean{ArakelovGeometry.AdelicStructure.logHeight_comp_aut,ArakelovGeometry.AdelicStructure.affineHeight_comp_aut}
\leanok
Let $S$ be any adelic structure, let $E/K$ be finite and $L/K$
algebraic, and let $i:E\hookrightarrow L$ be a $K$-embedding.
For $\sigma\in\operatorname{Aut}_K(E)$ and a tuple
$a\in E^I$, where $I$ is finite and nonempty, one has
\[
 h_{S\otimes_K L}\bigl((i(\sigma(a_j)))_{j\in I}\bigr)
 =h_{S\otimes_K L}\bigl((i(a_j))_{j\in I}\bigr).
\]
In particular, for every $\alpha\in E$,
\[
 h_{S\otimes_K L}^{\mathrm{aff}}(i(\sigma(\alpha)))
 =h_{S\otimes_K L}^{\mathrm{aff}}(i(\alpha)).
\]
For proper $S$, a finite normal subextension
$E\subseteq\overline K$, and its inclusion $i:E\hookrightarrow\overline K$,
this specializes to Lemma 3.5.4:
\[
 h_S([1:\sigma(\alpha)])=h_S([1:\alpha]).
\]
\end{lemma}
\begin{proof}
Apply Lemma~\ref{lem:height-finite-embedding} to $a$, first with
$i$ and then with $i\circ\sigma$. Both heights equal
$h_{S\otimes_K E}(a)$. Taking $a=(1,\alpha)$, and using that
field embeddings and automorphisms send $1$ to $1$, gives the affine
equality. When $S$ is proper and $L=\overline K$, these affine
heights equal the heights of the indicated projective representatives.

The finite-embedding measure identity used here is formulated with
separable-degree weights and holds for arbitrary finite extensions.
Thus the argument includes inseparable extensions directly; neither
separability nor a transitive action on the places is required. In
fact the tuple and affine assertions above require neither properness
nor the normality hypothesis in the stated source result.
\end{proof}

\begin{definition}[Polynomial height]
\label{def:polynomial-height}
\source{chen-moriwaki-adelic:page-0216}
\uses{def:absolute-projective-height,lem:local-height-integrable,prop:height-scalar-invariant}
\notready
For a nonzero polynomial $F=\sum_{j=0}^d a_jX^j\in\overline K[X]$
of degree $d$, put
\[
 \|F\|_\chi=\max_{0\leq j\leq d}|a_j|_\chi,
 \qquad
 h_S(F)=h_S([a_d:\cdots:a_0])
 =\int_{\Omega_{\overline K}}\log\|F\|_\chi\,d\nu_{\overline K}(\chi).
\]
The norm $\|F\|_\chi$ is positive and its logarithm is integrable.
Adding zero coefficients or permuting the coefficient list leaves
the height unchanged, and $h_S(cF)=h_S(F)$ for every $c\ne0$.
\end{definition}
\begin{proof}
At least one coefficient is nonzero, so its strictly positive absolute
value bounds the maximum from below. Integrability and the integral
formula follow from the corresponding assertions for nonzero tuples.
Adding zero entries does not change a maximum of nonnegative numbers,
and a permutation does not change its set of values. Multiplying every
coefficient by $c$ does not change the projective height, by scalar
invariance.
\end{proof}

\begin{lemma}[Heights of linear factors]
\label{lem:polynomial-height-linear}
\uses{def:polynomial-height,def:affine-height,lem:affine-height-power}
\notready
For $\alpha\in\overline K$, one has
$h_S(X-\alpha)=h_S([1:\alpha])$.
\end{lemma}
\begin{proof}
The coefficient tuple of $X-\alpha$ is $(1,-\alpha)$, whose local
maximum equals $\max(1,|\alpha|_\chi)$ because
$|-\alpha|_\chi=|\alpha|_\chi$. Their logarithmic integrands, and
hence their heights, are equal.
\end{proof}

\begin{lemma}[Height bounds for monic coefficients]
\label{lem:polynomial-height-monic-coefficient}
\uses{def:polynomial-height,prop:height-coordinate-bound}
\notready
If $F=X^d+a_{d-1}X^{d-1}+\cdots+a_0\in\overline K[X]$ is monic,
then $h_S([1:a_j])\leq h_S(F)$ for every $0\leq j<d$.
\end{lemma}
\begin{proof}
In the coefficient tuple, the leading coordinate is $1$. The affine
coordinate bound applied with this coordinate as denominator gives
the asserted inequality for each $a_j$.
\end{proof}

\begin{definition}[The integrated archimedean exponent]
\label{def:height-archimedean-constant}
\uses{lem:archimedean-bound}
\notready
Let $\kappa_S$ be the archimedean exponent, extended by zero on
the nonarchimedean locus. Define the nonnegative finite real number
\[
 A_S=\int_\Omega\kappa_S(\omega)\,d\nu(\omega).
\]
\end{definition}
\begin{proof}
The exponent is a nonnegative integrable real function by
Lemma~\ref{lem:archimedean-bound}; its integral is consequently
a finite nonnegative real number.
\end{proof}

\begin{lemma}[The exponent in an algebraic extension]
\label{lem:height-archimedean-extension}
\uses{def:height-archimedean-constant,lem:nonarchimedean-restriction,lem:archimedean-complex-embedding,prop:algebraic-base-pullback}
\notready
For an algebraic extension $L/K$ and $x\in\Omega_L$,
\[
 \kappa_{S\otimes_K L}(x)=\kappa_S(\pi_{L/K}(x)).
\]
In particular $A_{S\otimes_K L}=A_S$.
\end{lemma}
\begin{proof}
Restriction preserves the nonarchimedean property in both directions,
so either both exponents vanish or both places are archimedean. In
the latter case the characteristic is zero. The two absolute values
agree on $2$, and their exponents $s,t$ satisfy $2^s=2^t$. Taking
logarithms and using $\log2>0$ yields $s=t$. The equality of integrals
now follows from pullback of the integrable base function $\kappa_S$.
\end{proof}

\begin{lemma}[Local norm of a polynomial product]
\label{lem:polynomial-local-product-bound}
\uses{def:polynomial-height,lem:height-archimedean-extension,lem:archimedean-complex-embedding,lem:powered-complex-sum-bound}
\notready
Let $F,G\in\overline K[X]\setminus\{0\}$ and set
$m=\min(\deg F+1,\deg G+1)\geq1$. For
$\chi\in\Omega_{\overline K}$ above $\omega\in\Omega$,
\[
 \|FG\|_\chi\leq m^{\kappa_S(\omega)}\|F\|_\chi\|G\|_\chi.
\]
\end{lemma}
\begin{proof}
Write $F=\sum_{i=0}^d a_iX^i$ and $G=\sum_{j=0}^e b_jX^j$.
The coefficient of $X^k$ in their product is the sum of the terms
$a_ib_j$ with $i+j=k$, $0\leq i\leq d$, and $0\leq j\leq e$.
There are at most $m$ such terms, since either index determines
the other. Every term has absolute value at most
$M=\|F\|_\chi\|G\|_\chi$.

At a nonarchimedean place, the ultrametric inequality bounds each
coefficient by $M$; here the exponent is zero. At an archimedean
place, write $|b|_\chi=|\tau(b)|^s$ for a complex embedding $\tau$,
where $s=\kappa_S(\omega)>0$. If the coefficient sum has $r>0$
terms, the powered finite-sum bound gives absolute value at most
$r^sM\leq m^sM$. An empty sum is zero and satisfies the same
inequality. Taking the maximum over all coefficients proves the claim.
\end{proof}

\begin{lemma}[Polynomial height bound with finite error]
\label{lem:polynomial-height-product-sharp}
\uses{def:polynomial-height,def:height-archimedean-constant,lem:height-archimedean-extension,lem:polynomial-local-product-bound}
\notready
For nonzero $F,G\in\overline K[X]$, with
$m=\min(\deg F+1,\deg G+1)$, one has the real inequality
\[
 h_S(FG)\leq h_S(F)+h_S(G)+(\log m)A_S.
\]
\end{lemma}
\begin{proof}
All three polynomial norms are positive. Taking logarithms in the
local product estimate gives
\[
 \log\|FG\|_\chi\leq
 \log\|F\|_\chi+\log\|G\|_\chi
 +(\log m)\kappa_{S\otimes_K\overline K}(\chi).
\]
Each term is integrable: the first three are local heights of
coefficient tuples and the last is a constant multiple of the
integrable exponent. Integrate the inequality, use additivity of
integrals, and replace the integral of the extended exponent by $A_S$.
\end{proof}

\begin{lemma}[Polynomial height bound using archimedean mass]
\label{lem:polynomial-height-product}
\dcref{Lemma 3.5.5}
\source{chen-moriwaki-adelic:page-0217}
\uses{lem:polynomial-height-product-sharp,lem:archimedean-bound}
\notready
Let $F,G\in\overline K[X]\setminus\{0\}$ and
$m=\min(\deg F+1,\deg G+1)$. Then
\[
 h_S(FG)\leq h_S(F)+h_S(G)+(\log m)\nu(\Omega_\infty).
\]
When $\nu(\Omega_\infty)=\infty$, the right side is interpreted
in $[0,\infty]$, with $0\cdot\infty=0$. When this measure is finite,
the displayed formula is an inequality of real numbers.
\end{lemma}
\begin{proof}
If $m=1$, its logarithm is zero and the sharper estimate is exactly
the displayed inequality. Suppose $m>1$. If the archimedean mass is
infinite, its product with $\log m>0$ is infinite and the inequality
holds. If the mass is finite, the pointwise inequality
$0\leq\kappa_S\leq\mathbf1_{\Omega_\infty}$ implies
$A_S\leq\nu(\Omega_\infty)$. Multiply by $\log m\geq0$ and use
Lemma~\ref{lem:polynomial-height-product-sharp}.
\end{proof}

\begin{remark}[Finite error without a finite archimedean mass]
\label{rem:northcott-archimedean-error}
\source{chen-moriwaki-adelic:page-0217}
\uses{lem:polynomial-height-product,lem:polynomial-height-product-sharp}
\notready
The archimedean-mass bound is the formula in Lemma 3.5.5 of
\cite{chen-moriwaki}. Integrability of $\kappa_S$ does not impose
finiteness of $\nu(\Omega_\infty)$, since its positive values need
not be bounded away from zero. The proof of Northcott finiteness
below therefore uses the stronger finite error $(\log m)A_S$.
\end{remark}

\begin{lemma}[Roots of an irreducible polynomial are conjugate]
\label{lem:minimal-polynomial-root-conjugacy}
\uses{lem:algebraic-embedding-extension,lem:height-conjugate}
\notready
Let $F\in K[X]$ be monic and irreducible of degree $d\geq1$, let
$\alpha\in\overline K$ be a root, and write
$F=\prod_{j=1}^d(X-\alpha_j)$ in $\overline K[X]$, with roots
listed with multiplicity. Then
$h_S([1:\alpha_j])=h_S([1:\alpha])$ for every $j$.
This assertion includes inseparable irreducible polynomials.
\end{lemma}
\begin{proof}
Let $E$ be the subfield of $\overline K$ generated over $K$ by
the roots of $F$. It is finite over $K$, since it is generated by
finitely many algebraic elements. Each $K$-embedding of $E$ into
$\overline K$ maps the finite set of distinct roots of $F$
injectively into itself and therefore permutes that set. It follows
that the image is $E$, so $E/K$ is normal.

For a root $\beta$ of $F$, the assignments $X\mapsto\alpha$ and
$X\mapsto\beta$ identify both $K(\alpha)$ and $K(\beta)$ with
$K[X]/(F)$. Consequently there is a $K$-embedding
$K(\alpha)\to E$ sending $\alpha$ to $\beta$. Extend it to an
embedding $E\to\overline K$. The preceding permutation argument
makes this extension an automorphism of $E$. Apply
Lemma~\ref{lem:height-conjugate} to obtain equality of the two
heights. Every entry in the root list is a root, so the equality
holds also for repeated entries. No count of distinct roots or
separability assumption was used.
\end{proof}

\begin{lemma}[Height of a product of linear factors]
\label{lem:polynomial-height-split-bound}
\uses{lem:polynomial-height-linear,lem:polynomial-height-product-sharp,def:height-archimedean-constant}
\notready
If $d\geq1$ and $\alpha_1,\ldots,\alpha_d\in\overline K$, then
\[
 h_S\left(\prod_{j=1}^d(X-\alpha_j)\right)
 \leq\sum_{j=1}^d h_S([1:\alpha_j])+(d-1)(\log2)A_S.
\]
\end{lemma}
\begin{proof}
For $d=1$ this is the linear-factor equality. If the bound holds
for $d=r\geq1$, its product is monic of degree $r$. Multiplying
by $X-\alpha_{r+1}$ gives $m=\min(r+1,2)=2$ in the sharp
product bound. Add $h_S(X-\alpha_{r+1})+(\log2)A_S$ to the
inductive bound and use the linear-factor equality. The sum of the
two error coefficients is $(r-1)+1=r$, which gives the assertion
for $r+1$. Induction proves every case.
\end{proof}

\begin{lemma}[Height of the minimal polynomial]
\label{lem:minimal-polynomial-height}
\source{chen-moriwaki-adelic:page-0217}
\uses{lem:minimal-polynomial-root-conjugacy,lem:polynomial-height-split-bound,lem:polynomial-height-monic-coefficient,prop:height-algebraic-extension}
\notready
Let $\alpha\in\overline K$, let $d=[K(\alpha):K]$, and let
$F=X^d+a_{d-1}X^{d-1}+\cdots+a_0\in K[X]$ be its monic
minimal polynomial. Then
\[
 h_S(F)\leq d\,h_S([1:\alpha])+(d-1)(\log2)A_S,
 \qquad
 h_S^{\mathrm{aff}}(a_j)\leq h_S(F)\quad(0\leq j<d).
\]
\end{lemma}
\begin{proof}
The degree $d$ is positive, and $F$ factors over $\overline K$
as a product of $d$ linear factors, counting multiplicities.
Every root has the same height as $\alpha$ by
Lemma~\ref{lem:minimal-polynomial-root-conjugacy}. Substitute
these equal heights in the split-product bound. Since $F$ is monic,
the monic-coefficient bound gives the second inequality with
$h_S([1:a_j])$ on the left. The height of a base-field coefficient
agrees with $h_S^{\mathrm{aff}}(a_j)$ by algebraic-extension
invariance. This is the minimal-polynomial estimate in the
proof of Claim 3.5.6, strengthened by using $A_S$ in its error term.
\end{proof}

\begin{lemma}[Northcott finiteness for algebraic elements]
\label{lem:northcott-algebraic-elements}
\dcref{Claim 3.5.6}
\source{chen-moriwaki-adelic:page-0217}
\uses{def:northcott-property,lem:affine-height-power,lem:minimal-polynomial-height,def:height-archimedean-constant}
\notready
Assume $S$ has the Northcott property. For every $C,\delta\in\mathbb R$,
the set
\[
 \{\alpha\in\overline K:
       h_S([1:\alpha])\leq C,\ [K(\alpha):K]\leq\delta\}
\]
is finite.
\end{lemma}
\begin{proof}
An affine height is nonnegative and a simple field-extension degree
is a positive integer. The set is therefore empty if $C<0$ or
$\delta<1$. Otherwise let $D=\lfloor\delta\rfloor\geq1$ and put
\[
 B=\delta C+(\delta-1)(\log2)A_S\geq0.
\]
This is a finite real number. For an element in the stated set, its
minimal polynomial has degree $1\leq d\leq D$, and each nonleading
coefficient $a_j$ satisfies
\[
 h_S^{\mathrm{aff}}(a_j)\leq h_S(F)
 \leq d\,h_S([1:\alpha])+(d-1)(\log2)A_S\leq B.
\]
For the last inequality use $h_S([1:\alpha])\leq C$, $C\geq0$,
$d\leq\delta$, and $(\log2)A_S\geq0$.

By the Northcott property the set
$T=\{a\in K:h_S^{\mathrm{aff}}(a)\leq B\}$ is finite.
For each $d\in\{1,\ldots,D\}$ there are only finitely many
monic degree-$d$ polynomials with nonleading coefficients in $T$:
they are indexed by the finite set $T^d$. The union over the
finitely many possible degrees is finite. A nonzero polynomial of
degree $d$ has at most $d$ roots in a field: if a root $b$ exists,
factor out $X-b$; every other root is a root of the quotient, and
induction on the degree proves the bound, starting with a nonzero
constant. The union of the root sets of the finitely many possible
minimal polynomials is thus finite and contains the desired set.
\end{proof}

\begin{definition}[Field of definition of a projective point]
\label{def:projective-field-of-definition}
\source{chen-moriwaki-adelic:page-0216,chen-moriwaki-adelic:page-0217,chen-moriwaki-adelic:page-0218}
\notready
Let $x=[a_0:\cdots:a_n]\in\mathbb P^n(\overline K)$ and
choose $j$ with $a_j\ne0$. Define
\[
 K(x)=K(a_0/a_j,\ldots,a_n/a_j)\subseteq\overline K.
\]
This field does not depend on the representative or on $j$, and
it is finite over $K$. It is the smallest intermediate field $F$
of $\overline K/K$ over which $x$ has a projective representative.
\end{definition}
\begin{proof}
Replacing $a$ by $ca$ with $c\ne0$ leaves every ratio unchanged.
If $a_k\ne0$, then
$a_i/a_k=(a_i/a_j)/(a_k/a_j)$ belongs to the field defined using
$j$. This proves containment of the field defined using $k$ in
the field defined using $j$. Reversing $j$ and $k$ proves equality.
Each ratio is algebraic over $K$, and adjoining finitely many
algebraic elements gives a finite extension by the tower formula
for degrees.

The tuple $(a_i/a_j)_{i=0}^n$ has coordinate $1$ at $j$, is a
representative of $x$, and belongs to $K(x)^{n+1}$. Conversely,
if $b\in F^{n+1}\setminus\{0\}$ represents $x$, then $a=cb$
for a nonzero $c\in\overline K$. Thus $b_j\ne0$ and
$a_i/a_j=b_i/b_j\in F$ for every $i$. Hence $K(x)\subseteq F$,
proving the asserted minimality.
\end{proof}

\begin{lemma}[Degree and height bounds on projective charts]
\label{lem:projective-chart-height-degree}
\source{chen-moriwaki-adelic:page-0217,chen-moriwaki-adelic:page-0218}
\uses{def:projective-field-of-definition,prop:height-coordinate-bound,def:absolute-projective-height}
\notready
For $x=[a_0:\cdots:a_n]\in\mathbb P^n(\overline K)$ with
$a_j\ne0$, put $b_i=a_i/a_j$. Then
\[
 [K(b_i):K]\leq[K(x):K],\qquad
 h_S([1:b_i])\leq h_S(x)\quad(0\leq i\leq n).
\]
On the chart $a_j\ne0$, the normalized tuple $(b_i)_{i=0}^n$
determines $x$ uniquely.
\end{lemma}
\begin{proof}
There is a tower $K\subseteq K(b_i)\subseteq K(x)$ of finite
extensions. Its degrees multiply:
$[K(x):K]=[K(x):K(b_i)]\,[K(b_i):K]$.
The first factor is a positive integer, proving the degree bound.
Apply the affine coordinate-height bound to the representative $a$
over the proper curve $S\otimes_K\overline K$ to get the height
inequality. Finally $(b_i)$ represents $x$ and has coordinate $1$
at $j$. Two points with the same normalized tuple have the same
representative and are equal.
\end{proof}

\begin{theorem}[Northcott's theorem]
\label{thm:northcott}
\dcref{Theorem 3.5.3}
\source{chen-moriwaki-adelic:page-0216,chen-moriwaki-adelic:page-0217,chen-moriwaki-adelic:page-0218}
\uses{def:northcott-property,lem:northcott-algebraic-elements,def:projective-field-of-definition,lem:projective-chart-height-degree}
\notready
If $S$ has the Northcott property, then for every $n\in\mathbb N$
and every $C,\delta\in\mathbb R$, the set
\[
 \{x\in\mathbb P^n(\overline K):
       h_S(x)\leq C,\ [K(x):K]\leq\delta\}
\]
is finite.
\end{theorem}
\begin{proof}
By Lemma~\ref{lem:northcott-algebraic-elements}, the set
\[
 T=\{b\in\overline K:
       h_S([1:b])\leq C,\ [K(b):K]\leq\delta\}
\]
is finite. Fix $j\in\{0,\ldots,n\}$ and restrict to points
in the displayed set with nonzero $j$th coordinate. By
Lemma~\ref{lem:projective-chart-height-degree}, their normalized
tuples belong to $T^{n+1}$ and determine the points injectively.
This chart subset is finite because $T^{n+1}$ is finite. Every
projective point has some nonzero coordinate, so the full set is
the union of these $n+1$ finite chart subsets and is finite.
\end{proof}

\begin{corollary}[Northcott property under finite extension]
\label{cor:northcott-finite-extension}
\dcref{Corollary 3.5.7}
\source{chen-moriwaki-adelic:page-0218}
\uses{def:northcott-property,thm:northcott,def:projective-field-of-definition,lem:height-finite-embedding,lem:algebraic-embedding-extension,prop:finite-proper}
\notready
Assume $S$ has the Northcott property and let $E/K$ be a finite
extension. Then $S\otimes_K E$ has the Northcott property.
\end{corollary}
\begin{proof}
The extended adelic curve is proper. Choose a $K$-embedding
$i:E\hookrightarrow\overline K$, which exists by extension of
the inclusion $K\hookrightarrow\overline K$. Let $C\geq0$ and
$d=[E:K]$. For any $a\in E$, the field $K(i(a))$ is an
intermediate field of the finite degree-$d$ extension $i(E)/K$,
so the degree tower formula gives $[K(i(a)):K]\leq d$.
Finite-embedding compatibility also gives
\[
 h_{S\otimes_K E}^{\mathrm{aff}}(a)=h_S([1:i(a)]).
\]
For $x=[1:i(a)]$, the ratio definition gives $K(x)=K(i(a))$.
Consequently $a\mapsto[1:i(a)]$ maps the elements of $E$ with
affine height at most $C$ into the set of projective points of
height at most $C$ and degree at most $d$ in
$\mathbb P^1(\overline K)$. This map is injective: proportional
representatives whose first coordinates are both $1$ have scalar
factor $1$, and $i$ is injective. The target is finite by
Theorem~\ref{thm:northcott}; its subset and the injectively mapped
source set are finite. Since this holds for every $C\geq0$, the
extended curve has the Northcott property.
\end{proof}
